\setlength{\baselineskip}{20pt}
\chapter{相关工作}
\label{cha:chap2}

\section{生成式扩散模型}
扩散模型通过前向加噪与反向去噪过程学习数据分布，近年来在图像与多模态生成领域取得了显著进展。以 DDPM 为代表的方法将采样过程离散化为多步马尔可夫链，具有良好的训练稳定性与生成质量。随后，DDIM、Progressive Distillation、DPM-Solver 等工作进一步从采样轨迹压缩和数值求解角度降低了推理开销。上述研究表明，扩散框架不仅具备高质量建模能力，也具备较强的噪声估计能力，这为其向判别任务迁移提供了理论基础。

\section{判别式表征学习}
判别任务以学习条件分布 $P(y|x)$ 为目标，高质量表征是分类、检索、检测、分割等任务性能的共同基础。传统方法从手工特征与度量学习出发，深度学习阶段逐步形成以 CNN 与 Transformer 为核心的主流框架。行人重识别领域尤为依赖特征判别性，研究从 ResNet 系列方法扩展到 TransReID、CLIP-ReID 等新范式。尽管性能持续提升，复杂场景下的噪声累积与分布偏移问题仍普遍存在。

\section{扩散模型在判别任务中的应用}
已有研究尝试将扩散思想用于检测与分割等特定任务，例如将检测框逐步去噪或利用扩散注意力图构建零样本分割结果。这些工作验证了扩散思想在判别任务中的潜力，但大多依赖任务特定数据结构，难以直接迁移到通用表征学习框架。同时，多步迭代去噪常引入额外推理时延，与实际部署需求存在矛盾。

\section{知识蒸馏相关方法}
知识蒸馏通过教师网络向学生网络传递软监督信号，可在参数规模受限条件下提升学生模型表达能力。扩散模型领域的渐进式蒸馏表明，多步去噪轨迹可以通过蒸馏压缩到更少步数并保持性能。该思路对本文具有直接启发：在不改变部署复杂度的前提下，利用教师模型提升轻量去噪模块的噪声估计精度，并通过多步信息压缩增强单步推理能力。

\section{本章小结}
本章从生成式扩散模型、判别式表征学习、扩散判别迁移与知识蒸馏四个维度总结了相关研究现状。现有方法在“通用性、效率与去噪强度”之间仍存在明显缺口。下一章将提出基础方法 DenoiseRep，并在方法章内直接给出对应实验验证。
